\section{The algebra determined by one three-plane}\label{sec:algebra}

We first isolate the algebraic implication of the analytic theorem.
Let $X$ be a projective primitive symplectic variety of complex
dimension $2n$, with $b_2(X)\ge5$ and an injection
$H^2(X;\QQ)\to\IH^2(X;\QQ)$. Set
\[
 H=H^2(X;\CC),\qquad V_X=\IH^\bullet(X;\CC),\qquad q=q_X.
\]
The letter $H$ in this section denotes cohomology; in the quotient
constructions below it will denote a compact stabiliser. We retain
the grading operator $h|_{\IH^j}=j-2n$ and the notation $L_x$,
$\Lambda_x$ from the introduction.

\begin{theorem}\label{thm:criterion}
Suppose that three real Lefschetz classes $x_1,x_2,x_3$ satisfy
$q(x_a,x_b)=c\delta_{ab}$ with $c>0$. Put
$L_a=L_{x_a}$, $\Lambda_a=\Lambda_{x_a}$ and
$K_{ab}=[L_a,\Lambda_b]$ for $a\ne b$. Assume
\begin{equation}\label{eq:V}
 [\Lambda_a,\Lambda_b]=0,\qquad K_{ab}=-K_{ba},\qquad
 [K_{ab},L_c]=2(\delta_{bc}L_a-\delta_{ac}L_b).
\end{equation}
Then the LLV algebra of $X$ is
$\mathfrak{so}((H,q)\oplus U)$, with graded pieces
\[
 \mathfrak g_2=L(H),\qquad
 \mathfrak g_0\cong\mathfrak{so}(H,q)\oplus\CC h,
 \qquad \mathfrak g_{-2}\cong H.
\]
The isomorphism is defined over $\QQ$. Its real form is
$\mathfrak{so}(4,b_2(X)-2)$.
\end{theorem}

\subsection{The geometric inputs to the criterion}
We use three established facts. Locally trivial deformations give
parallel transport on ordinary and intersection cohomology, compatible
with their module structure. Indeed these families are locally
real-analytic fibre bundles \cite[Proposition~5.1]{av21}, and
intersection cohomology and its ordinary-cohomology action are
topological invariants \cite{gm83}. Consequently their monodromy
group satisfies
\begin{equation}\label{eq:mon}
 \gamma L_x\gamma^{-1}=L_{\gamma x}.
\end{equation}
The image on $H^2$ has finite index in the integral orthogonal group
\cite[Theorem~1.1]{bl22}; for the moduli spaces under consideration
see also \cite{opr24}. Borel density implies that its complex Zariski
closure contains $\SO(H,q)$.

Next, the Lefschetz locus is
\begin{equation}\label{eq:lefschetz-locus}
 H_{\mathrm{Lef}}=\{x\in H:q(x)\ne0\}.
\end{equation}
For rational classes this is \cite[Proposition~2.11]{tighe22}.
To pass to complex classes, the non-Lefschetz locus is closed,
conical and monodromy invariant, since it is defined by determinants
of the maps $L_x^j$. It is therefore $\SO(H,q)$-invariant. The
nonisotropic vectors form one orbit up to scaling, and this orbit
contains an ample class. The nonzero isotropic vectors form one
orbit. This orbit contains a rational point by Meyer's theorem,
because the rational form is indefinite and has rank at least five.
The rational statement proves \eqref{eq:lefschetz-locus}.

Finally $\mathfrak g$ is semisimple. This is the polarised Lefschetz
module criterion of \cite[Proposition~1.6]{ll97}: Saito's pure Hodge
structures on $\IH$, the Hodge-compatible action of $H^2$, and the
Hodge--Riemann relations for an ample class supply its hypotheses
\cite{saito88,saito90,dcm05}. The intersection pairing is taken with
the degree signs of \cite[\S1.9]{ll97}, so that the raising
operators are skew. These facts concern intersection cohomology
algebraically and make no assertion about types of harmonic forms.

\subsection{Extending the relations}
For a Lefschetz class the equations defining $\Lambda_x$ have a
unique solution. Solving this finite linear system shows that
$x\mapsto\Lambda_x$ is a regular rational map on
$H_{\mathrm{Lef}}$ \cite[\S1.1]{ll97}. Uniqueness also gives
$\Lambda_{tx}=t^{-1}\Lambda_x$ and the conjugacy analogue of
\eqref{eq:mon}.

\begin{lemma}\label{lem:density}
Let $N\le3$ and let $\mathcal D\subset H^N$ be a dense invariant
Zariski open set. A relatively closed monodromy-invariant subset
of $\mathcal D$ which contains
$\mathcal D\cap\mathcal P_{\CC}^N$, where
$\mathcal P=\langle x_1,x_2,x_3\rangle$, is all of $\mathcal D$.
\end{lemma}
\begin{proof}
The stabiliser of a closed algebraic subset is Zariski closed, so the
subset is invariant under $\SO(H,q)$. Tuples with independent
entries and nondegenerate span are dense. Each such span is contained
in a nondegenerate three-plane. Witt's extension theorem carries
$\mathcal P_{\CC}$ to this plane by an orthogonal transformation.
A reflection in its orthogonal complement corrects the determinant
if necessary; that complement has positive dimension. Invariance
therefore places every such tuple in the closed subset, proving the
claim by density.
\end{proof}

For $x=\sum_at_ax_a\in\mathcal P_{\CC}$ the first two relations in
\eqref{eq:V} imply
\[
 [\sum_at_aL_a,\sum_at_a\Lambda_a]=(\sum_at_a^2)h.
\]
Uniqueness of the lowering operator gives
\begin{equation}\label{eq:plane}
 q(x)\Lambda_x=c\sum_at_a\Lambda_a.
\end{equation}
Thus lowering operators commute on this plane. Their commutator
vanishes everywhere on $H_{\mathrm{Lef}}^2$ by Lemma
\ref{lem:density}.

Define initially $A_x=q(x)\Lambda_x$ for $q(x)\ne0$. Equation
\eqref{eq:plane} gives additivity on pairs in $\mathcal P_{\CC}$
for which $x,y,x+y$ are nonisotropic. Lemma \ref{lem:density}
extends it to every such pair. Together with homogeneity this
extends $A$ to a linear map on $H$. For clarity, one may differentiate
$A(x+y)=A(x)+A(y)$ at generic pairs: $dA_{x+y}=dA_x=dA_y$,
so $dA$ is constant on its irreducible domain; homogeneity removes
the constant term.

Use the convention
$(x\wedge y)z=q(y,z)x-q(x,z)y$ for $\mathfrak{so}(H,q)$.
The last relation of \eqref{eq:V}, together with
$[h,L_z]=2L_z$, now gives on the three-plane
\[
 [[L_x,A_y],L_z]=2L_{q(x,y)z+(x\wedge y)z}.
\]
Both sides are regular for $q(y)\ne0$ before the linear extension
of $A$. Lemma \ref{lem:density} with three entries, followed by
linearity, therefore proves
\begin{equation}\label{eq:structure}
 [A_x,A_y]=0,\qquad
 [[L_x,A_y],L_z]=2L_{q(x,y)z+(x\wedge y)z}
 \quad(x,y,z\in H).
\end{equation}

\subsection{Recovering the orthogonal algebra}
\begin{proof}[Proof of Theorem \ref{thm:criterion}]
The map $L:H\to\End(V_X)$ is injective, since $L_x1$ is the image
of $x$ in $\IH^2$. Let $\mathfrak k$ be generated by the brackets
$[L_x,A_y]$. Equation \eqref{eq:structure} defines a Lie
homomorphism $\pi:\mathfrak k\to\mathfrak{gl}(H)$ by
\[
 [T,L_z]=L_{\pi(T)z},\qquad
 \pi([L_x,A_y])=2q(x,y)\id+2x\wedge y.
\]
For $T\in\mathfrak k$, conjugation by $e^{sT}$ preserves $h$ and
takes $L_z$ to $L_{e^{s\pi(T)}z}$. It takes the corresponding
lowering operator to the corresponding lowering operator. Hence
\begin{equation}\label{eq:lowering-action}
 [T,\Lambda_z]=\frac{A_{\pi(T)z}}{q(z)}
 -\frac{2q(z,\pi(T)z)}{q(z)^2}A_z.
\end{equation}
It follows that $L(H)\oplus\mathfrak k\oplus A(H)$ is closed
under brackets and equals $\mathfrak g$. Its three summands have
degrees $2,0,-2$. Conjugation by the Weyl element of any Lefschetz
$\mathfrak{sl}_2$ exchanges the two extreme pieces, so $A$ is also
injective.

The image of $\pi$ is
$\mathfrak{so}(H,q)\oplus\CC\id$, and $\pi(h)=2\id$.
If $\pi(T)=0$, then $T$ commutes with all raising operators and,
by \eqref{eq:lowering-action}, all lowering operators. It is
central in the semisimple algebra $\mathfrak g$, and is zero.
Thus $\pi$ is an isomorphism. If $\pi(T)=a\id+B$ with
$B\in\mathfrak{so}(H,q)$, the remaining brackets are
\[
 [T,L_z]=L_{az+Bz},\qquad [T,A_z]=A_{-az+Bz},\qquad
 [L_x,A_y]=\pi^{-1}(2q(x,y)\id+2x\wedge y).
\]

Choose isotropic generators $u_+,u_-$ of $U$ with
$\widetilde q(u_+,u_-)=1$. All the brackets agree under
\begin{equation}\label{eq:orthogonal-map}
 L_x\longmapsto u_+\wedge x,\quad
 A_y\longmapsto-2u_-\wedge y,\quad
 \pi^{-1}(B)\longmapsto B,\quad
 h\longmapsto2u_+\wedge u_-.
\end{equation}
For example
$[u_+\wedge x,u_-\wedge y]=-x\wedge y-q(x,y)u_+\wedge u_-$.
The indicated elements exhaust the orthogonal algebra and have the
same graded dimensions. This proves the complex isomorphism.

The maps $L$, $A$ and $\pi$ are rational: their values on rational
nonisotropic classes are rational, and those classes are Zariski
dense. Formula \eqref{eq:orthogonal-map} is therefore defined over
$\QQ$. The real signature of $q\oplus U$ is $(4,b_2-2)$.
\end{proof}

\subsection{Application to the moduli space}
For the moment use Theorem \ref{thm:analytic}, proved below.
The bounded K\"ahler forms represent classes in $H^2(M;\RR)$.
Proposition \ref{prop:qF} proves
$q_M([\omega_a],[\omega_b])=c\delta_{ab}$ with $c>0$.
Under the harmonic representative isomorphism, $L_{[\omega_a]}$
is pointwise wedge product and its lowering operator is the
pointwise adjoint. Indeed these operators preserve harmonic forms
and constitute an $\mathfrak{sl}_2$-triple. Finite-dimensional
$\mathfrak{sl}_2$ theory gives hard Lefschetz, and uniqueness then
identifies the lowering operator. The pointwise hyperk\"ahler
relations are \eqref{eq:V}. Also
$H^2(M;\QQ)\cong\IH^2(M;\QQ)$ by
\cite[Proposition~2.17]{tighe22}. Theorem \ref{thm:criterion}
proves the algebra statement in Theorem \ref{thm:main}; its
dimension is $\binom{25}{2}=300$.

To identify the unit submodule, the negative part of $\mathfrak g$
annihilates $1$, the semisimple part of $\mathfrak g_0$ acts
trivially on it, and $h1=-2n$. Complete reducibility shows that a
module generated by this lowest-weight vector is the irreducible
$\VV_{(n)}$. The raising operators generate the same module, and
the usual description \cite[Proposition~2.14]{ll97} is
\begin{equation}\label{eq:verbitsky-component}
 \VV_{(n)}\cong
 \Sym^\bullet H^2/\langle x^{n+1}:q_M(x)=0\rangle.
\end{equation}
There can be only one such summand because $\IH^0(M)=\CC$.
The ideal has no element of degree at most $n$, giving the
injections and dimension bounds in Theorem \ref{thm:main}.
